\documentclass[11pt]{article}
\usepackage{amsmath,amssymb}

\title{Message-Passing Graph Neural Networks for Molecular Property Prediction}
\author{F. Iqbal \and N. Schreiber \and Y. Watanabe}
\date{}

\begin{document}
\maketitle

\begin{abstract}
We benchmark message-passing graph neural networks against fingerprint and
descriptor baselines for predicting eight quantum-chemical properties of
small organic molecules. With careful normalisation and an edge-feature
augmentation that encodes bond stereochemistry, a six-layer GNN reaches
chemical accuracy on six of the eight QM9 targets without any pre-training.
\end{abstract}

\section{Background}
A molecule is naturally a graph: atoms are nodes, bonds are edges, and many
physicochemical properties decompose locally over neighbourhoods. Classical
approaches hand-engineer Morgan fingerprints or RDKit descriptors and then
fit a gradient-boosted regressor. We replace the feature stage with a learned
message-passing layer and let the regressor share representation with the
encoder.

\section{Architecture}
Let $h_v^{(0)}$ be initial atom features (element, charge, hybridisation,
aromaticity). At layer $\ell$:
\begin{equation}
h_v^{(\ell+1)} = \phi\!\left( h_v^{(\ell)},\; \bigoplus_{u \in \mathcal{N}(v)}
  \psi(h_u^{(\ell)}, h_v^{(\ell)}, e_{uv}) \right),
\end{equation}
where $\bigoplus$ is a permutation-invariant aggregator (sum), $\psi$ is an
edge-conditioned MLP, and $\phi$ is a gated update with a residual term.
A graph-level read-out
\begin{equation}
h_G = \sigma\!\left(\sum_{v \in G} \alpha_v h_v^{(L)}\right),\quad
\alpha_v = \mathrm{softmax}_v(W_a h_v^{(L)})
\end{equation}
attention-pools node embeddings into a fixed-dimensional graph descriptor.

\section{Training}
We follow QM9's canonical 110k/10k/13k train/val/test split. Targets are
normalised per-property to zero mean and unit variance; predictions are
de-normalised at evaluation. Adam with $\eta = 5\times10^{-4}$ converges in
about 200 epochs on a single A100. Early stopping on validation MAE prevents
the typical late-stage overfitting on dipole moment.

\section{Results}
\begin{tabular}{lcc}
\hline
Target & RDKit + LightGBM & GNN (ours) \\
\hline
HOMO   & 0.092 eV  & 0.041 eV \\
LUMO   & 0.114 eV  & 0.038 eV \\
Dipole & 0.41 D    & 0.18 D \\
Energy & 1.8 kcal/mol & 0.6 kcal/mol \\
\hline
\end{tabular}

The GNN's largest win is on dipole moment, a property whose value depends on
charge distribution far from any single atom --- exactly the regime where
local fingerprints lose information.

\section{Ablations}
Removing edge features (bond order, stereochemistry) costs roughly $25\%$ of
the gain over baselines. Replacing sum-pool with mean-pool collapses
performance for the energy target; the absolute scale of an extensive quantity
genuinely depends on graph size, and mean-pooling discards exactly that.

\section{Discussion}
GNNs are sometimes oversold as universal molecular models. Our results suggest
the more measured claim: when the property is local, expressive aggregators
help; when it is extensive, the pooling operator must respect that scaling.
The right inductive bias is property-specific.

\bibliographystyle{plain}
\begin{thebibliography}{9}
\bibitem{gilmer2017neural} J. Gilmer et al. \emph{Neural Message Passing for Quantum Chemistry.} ICML, 2017.
\bibitem{ramakrishnan2014quantum} R. Ramakrishnan et al. \emph{Quantum chemistry structures and properties of 134 kilo molecules.} Sci. Data, 2014.
\end{thebibliography}

\end{document}
